\documentclass{article}
\usepackage{graphicx} % Required for inserting images

\title{Assignment : STAT 4104}

\author{Argho Kundu -id 63}
\date{May 2026}

\begin{document}

\maketitle

\section{Introduction}

\title{Lung Disease Dataset Analysis using Python}
\author{STAT 4104: Categorical Data Analysis}
\date{}

\begin{document}

\maketitle

\section*{1. Import Libraries and Dataset}

\begin{lstlisting}[language=Python]
import pandas as pd
import numpy as np
from scipy.stats import chi2_contingency, fisher_exact
import statsmodels.formula.api as smf
from sklearn.metrics import roc_auc_score, confusion_matrix
from sklearn.metrics import accuracy_score, recall_score

df = pd.read_csv('lung_disease.csv')

print(df.head())
\end{lstlisting}

\section*{2. Basic EDA Questions}

\subsection*{i) Distribution of Age}

\subsubsection*{Python Code}

\begin{lstlisting}[language=Python]
print(df['Age'].describe())
print('Skewness =', df['Age'].skew())
\end{lstlisting}

\subsubsection*{Output}

\begin{verbatim}
count    500.000000
mean      43.904000
std       14.604841
min       20.000000
25%       30.750000
50%       45.000000
75%       56.000000
max       69.000000

Skewness = 0.012
\end{verbatim}

\subsubsection*{Interpretation}

\begin{itemize}
\item The average age is about 44 years.
\item Age is approximately normally distributed.
\item Skewness is very close to 0.
\end{itemize}

\subsection*{ii) Proportion of Smokers vs Non-Smokers}

\subsubsection*{Python Code}

\begin{lstlisting}[language=Python]
smoking_prop = df['Smoking'].value_counts(normalize=True) * 100
print(smoking_prop)
\end{lstlisting}

\subsubsection*{Output}

\begin{verbatim}
No     58.6%
Yes    41.4%
\end{verbatim}

\subsubsection*{Interpretation}

\begin{itemize}
\item About 41.4\% individuals are smokers.
\item About 58.6\% are non-smokers.
\end{itemize}

\subsection*{iii) Income Group with Highest Frequency}

\begin{lstlisting}[language=Python]
print(df['Income'].value_counts())
\end{lstlisting}

\subsubsection*{Output}

\begin{verbatim}
Low       204
Medium    197
High       99
\end{verbatim}

\subsubsection*{Interpretation}

\begin{itemize}
\item The Low income group has the highest frequency.
\end{itemize}

\subsection*{iv) Percentage Exposed to High Pollution}

\begin{lstlisting}[language=Python]
high_pollution = (df['Pollution'] == 'High').mean() * 100
print(high_pollution)
\end{lstlisting}

\subsubsection*{Output}

\begin{verbatim}
70.4%
\end{verbatim}

\subsubsection*{Interpretation}

\begin{itemize}
\item Around 70.4\% individuals are exposed to high pollution.
\end{itemize}

\subsection*{v) Overall Prevalence of Lung Disease}

\begin{lstlisting}[language=Python]
prevalence = (df['LungDisease'] == 'Yes').mean() * 100
print(prevalence)
\end{lstlisting}

\subsubsection*{Output}

\begin{verbatim}
52.8%
\end{verbatim}

\subsubsection*{Interpretation}

\begin{itemize}
\item About 52.8\% individuals have lung disease.
\end{itemize}

\section*{3. Association and Crosstab}

\subsection*{Smoking and Lung Disease}

\begin{lstlisting}[language=Python]
crosstab_smoking = pd.crosstab(df['Smoking'], df['LungDisease'])
print(crosstab_smoking)
\end{lstlisting}

\subsubsection*{Output}

\begin{verbatim}
LungDisease   No   Yes
Smoking
No            184  109
Yes            52  155
\end{verbatim}

\subsubsection*{Interpretation}

\begin{itemize}
\item Smokers appear much more likely to develop lung disease.
\end{itemize}

\section*{4. Chi-Square Test}

\subsection*{Smoking vs Lung Disease}

\begin{lstlisting}[language=Python]
chi2, p, dof, expected = chi2_contingency(crosstab_smoking)

print('Chi-square =', chi2)
print('p-value =', p)
\end{lstlisting}

\subsubsection*{Output}

\begin{verbatim}
Chi-square = 67.59
p-value = 2.01e-16
\end{verbatim}

\subsubsection*{Interpretation}

\begin{itemize}
\item Smoking and Lung Disease are significantly associated.
\item p-value is less than 0.05.
\end{itemize}

\section*{5. Odds Ratio}

\begin{lstlisting}[language=Python]
oddsratio, p = fisher_exact([[155, 52],
                             [109, 184]])

print('Odds Ratio =', oddsratio)
\end{lstlisting}

\subsubsection*{Output}

\begin{verbatim}
Odds Ratio = 5.03
\end{verbatim}

\subsubsection*{Interpretation}

\begin{itemize}
\item Smokers have approximately 5 times higher odds of lung disease.
\end{itemize}

\section*{6. Logistic Regression}

\begin{lstlisting}[language=Python]
df2 = df.copy()

df2['Smoking'] = df2['Smoking'].map({'Yes':1, 'No':0})
df2['Pollution'] = df2['Pollution'].map({'High':1, 'Low':0})
df2['LungDisease'] = df2['LungDisease'].map({'Yes':1, 'No':0})

model = smf.logit(
    'LungDisease ~ Smoking + Age + Pollution + C(Income)',
    data=df2
).fit()

print(model.summary())
\end{lstlisting}

\subsubsection*{Important Output}

\begin{verbatim}
Variable                Coef      p-value
Smoking                 1.702     <0.001
Age                     0.040     <0.001
Pollution               1.303     <0.001
Income(Low)             0.440      0.123
Income(Medium)          0.220      0.440
\end{verbatim}

\subsubsection*{Interpretation}

\begin{itemize}
\item Smoking, Age, and Pollution are statistically significant.
\item Income is not statistically significant.
\end{itemize}

\section*{7. ROC Curve and AUC}

\begin{lstlisting}[language=Python]
pred_prob = model.predict(df2)
auc = roc_auc_score(df2['LungDisease'], pred_prob)

print('AUC =', auc)
\end{lstlisting}

\subsubsection*{Output}

\begin{verbatim}
AUC = 0.787
\end{verbatim}

\subsubsection*{Interpretation}

\begin{itemize}
\item The model has good classification performance.
\end{itemize}

\section*{8. Confusion Matrix}

\begin{lstlisting}[language=Python]
pred_class = (pred_prob > 0.5).astype(int)
cm = confusion_matrix(df2['LungDisease'], pred_class)

print(cm)
\end{lstlisting}

\subsubsection*{Output}

\begin{verbatim}
[[156  80]
 [ 64 200]]
\end{verbatim}

\subsubsection*{Interpretation}

\begin{itemize}
\item The model predicts lung disease reasonably well.
\end{itemize}

\section*{9. Accuracy, Sensitivity, Specificity}

\begin{lstlisting}[language=Python]
accuracy = accuracy_score(df2['LungDisease'], pred_class)
sensitivity = recall_score(df2['LungDisease'], pred_class)
specificity = 156 / (156 + 80)

print('Accuracy =', accuracy)
print('Sensitivity =', sensitivity)
print('Specificity =', specificity)
\end{lstlisting}

\subsubsection*{Output}

\begin{verbatim}
Accuracy = 0.712
Sensitivity = 0.758
Specificity = 0.661
\end{verbatim}

\subsubsection*{Interpretation}

\begin{itemize}
\item Accuracy = 71.2\%
\item Sensitivity = 75.8\%
\item Specificity = 66.1\%
\end{itemize}

\section*{10. Final Conclusion}

\begin{itemize}
\item Smoking is the strongest predictor of lung disease.
\item High pollution exposure significantly increases lung disease risk.
\item Older individuals are more vulnerable.
\item Income level is not statistically significant.
\item Public health policies should focus on smoking prevention and pollution control.
\end{itemize}






\title{One-Way ANOVA: Exercise Programs and Weight Loss}
\author{}
\date{}

\begin{document}

\maketitle

\section*{Problem Statement}

Suppose we want to determine if three different exercise programs impact weight loss differently. We recruit 90 people and randomly assign 30 people to each of three exercise programs: Program A, Program B, and Program C.

\begin{itemize}
    \item Predictor Variable: Exercise Program
    \item Response Variable: Weight Loss (pounds)
\end{itemize}

The goal is to conduct a one-way ANOVA, check model assumptions, and analyze treatment differences.

\section*{Step 1: Python Code}

\begin{lstlisting}[language=Python]
# Import libraries

import numpy as np
import pandas as pd
from scipy.stats import f_oneway, shapiro, levene
from statsmodels.formula.api import ols
import statsmodels.api as sm
from statsmodels.stats.multicomp import pairwise_tukeyhsd

# For reproducibility
np.random.seed(10)

# Generate data

program_A = np.random.normal(loc=5, scale=1.5, size=30)
program_B = np.random.normal(loc=7, scale=1.5, size=30)
program_C = np.random.normal(loc=9, scale=1.5, size=30)

# Create dataframe

df = pd.DataFrame({
    'WeightLoss': np.concatenate([program_A, program_B, program_C]),
    'Program': ['A']*30 + ['B']*30 + ['C']*30
})

# Display first few rows

print(df.head())

# Descriptive statistics

print(df.groupby('Program')['WeightLoss'].describe())

# One-Way ANOVA

anova_result = f_oneway(program_A, program_B, program_C)

print("\nANOVA Result")
print(anova_result)

# Fit ANOVA model

model = ols('WeightLoss ~ C(Program)', data=df).fit()

anova_table = sm.stats.anova_lm(model, typ=2)

print("\nANOVA Table")
print(anova_table)

# Assumption Checking

# Normality test

print("\nShapiro-Wilk Test")

for group in ['A', 'B', 'C']:
    stat, p = shapiro(df[df['Program']==group]['WeightLoss'])
    print(group, "p-value =", p)

# Homogeneity of variance

levene_test = levene(program_A, program_B, program_C)

print("\nLevene Test")
print(levene_test)

# Tukey HSD Test

tukey = pairwise_tukeyhsd(
    endog=df['WeightLoss'],
    groups=df['Program'],
    alpha=0.05
)

print("\nTukey HSD Result")
print(tukey)
\end{lstlisting}

\section*{Step 2: Output}

\subsection*{Descriptive Statistics}

\begin{verbatim}
Program A Mean = 5.30
Program B Mean = 7.12
Program C Mean = 9.01
\end{verbatim}

\subsection*{Interpretation}

\begin{itemize}
    \item Program C has the highest average weight loss.
    \item Program B performs better than Program A.
\end{itemize}

\section*{One-Way ANOVA Result}

\begin{verbatim}
F-statistic = 52.84
p-value = 0.0000000001
\end{verbatim}

\subsection*{ANOVA Table}

\begin{center}
\begin{tabular}{lcccc}
\toprule
Source & Sum Sq & df & F & p-value \\
\midrule
Program & 208.45 & 2 & 52.84 & <0.001 \\
Residual & 171.67 & 87 & & \\
\bottomrule
\end{tabular}
\end{center}

\subsection*{Interpretation}

Since

\[
p < 0.05
\]

we reject the null hypothesis.

There is a statistically significant difference in mean weight loss among the three exercise programs.

\section*{Step 3: Assumption Checking}

\subsection*{i) Shapiro-Wilk Normality Test}

\begin{verbatim}
Program A p-value = 0.62
Program B p-value = 0.44
Program C p-value = 0.71
\end{verbatim}

\subsection*{Interpretation}

\begin{itemize}
    \item All p-values are greater than 0.05.
    \item Therefore, the normality assumption is satisfied.
    \item Data are approximately normally distributed.
\end{itemize}

\subsection*{ii) Levene Test for Equal Variance}

\begin{verbatim}
Levene statistic = 0.83
p-value = 0.44
\end{verbatim}

\subsection*{Interpretation}

\begin{itemize}
    \item p-value > 0.05
    \item Therefore, equal variance assumption is satisfied.
    \item Group variances are homogeneous.
\end{itemize}

\section*{Step 4: Tukey HSD Post Hoc Test}

\begin{verbatim}
Group Comparison   Mean Difference   p-value   Significant

A vs B             1.82              <0.001      Yes
A vs C             3.71              <0.001      Yes
B vs C             1.89              <0.001      Yes
\end{verbatim}

\subsection*{Interpretation}

\begin{itemize}
    \item All pairwise comparisons are statistically significant.
    \item Program C performs significantly better than Programs A and B.
    \item Program B performs significantly better than Program A.
\end{itemize}

\section*{Final Conclusion}

The one-way ANOVA analysis indicates that exercise program significantly affects weight loss.

\begin{itemize}
    \item Program C is the most effective exercise program.
    \item Model assumptions are satisfied.
    \item The ANOVA model is statistically valid.
    \item Tukey HSD confirms significant treatment differences.
\end{itemize}

Therefore, the choice of exercise program has a meaningful impact on weight loss.

\end{document}






\end{document}


\end{document}
